
\documentclass[11pt]{article}
\usepackage[margin=1in]{geometry}
\usepackage{amsmath,amssymb,amsthm}
\usepackage{hyperref}
\usepackage{enumitem}

\newtheorem{definition}{Definition}
\newtheorem{lemma}{Lemma}
\newtheorem{proposition}{Proposition}
\newtheorem{theorem}{Theorem}
\newtheorem{remark}{Remark}

\title{Linearized Wilson Action Hessian at Vacuum and Its Spectrum (Finite Volume)}
\date{}

\begin{document}
\maketitle

\paragraph{Dependency.}
This document uses the configuration-space notation and the vacuum gauge projector $P_0$ from \emph{Document~01}.

\section{Wilson action}

Fix $\beta>0$.  For $U\in\mathcal M_L$, define the Wilson action
\[
S_W(U) := \beta\sum_{x\in\Lambda_L}\sum_{0\le \mu<\nu\le 3}\left(1-\frac13\mathrm{Re}\,\mathrm{Tr}\,U_{\mu\nu}(x)\right),
\]
where the plaquette variable is
\[
U_{\mu\nu}(x):=U_\mu(x)\,U_\nu(x+\hat\mu)\,U_\mu(x+\hat\nu)^{-1}\,U_\nu(x)^{-1}.
\]

\section{Right-invariant small-field coordinates}

Work in right-trivialized coordinates near the vacuum $U\equiv I$:
\[
U_\mu(x)=\exp(A_\mu(x)),\qquad A_\mu(x)\in \mathfrak{su}(3),\qquad \|A\|\ll 1.
\]

\begin{definition}[Discrete forward difference]
For a lattice field $f:\Lambda_L\to V$ (any vector space $V$), define
\[
(\Delta_\mu f)(x):=f(x+\hat\mu)-f(x).
\]
\end{definition}

\begin{definition}[Discrete exterior derivative on $1$-forms]
For a $1$-form $A=\{A_\mu(x)\}$, define the $2$-form $dA$ by
\[
(dA)_{\mu\nu}(x):=(\Delta_\mu A_\nu)(x)-(\Delta_\nu A_\mu)(x),\qquad \mu<\nu.
\]
\end{definition}

\section{Quadratic expansion}

\begin{lemma}[Class-function expansion at identity]\label{lem:class}
Let $X\in\mathfrak{su}(3)$ be small and set $V=\exp(X)\in SU(3)$.
Then
\[
1-\frac13\mathrm{Re}\,\mathrm{Tr}\,V = -\frac16\,\mathrm{Tr}(X^2) + O(\|X\|^3).
\]
With the inner product $\langle A,B\rangle=-\mathrm{Tr}(AB)$, this can be written as
\[
1-\frac13\mathrm{Re}\,\mathrm{Tr}\,\exp(X) = \frac16\,\langle X,X\rangle + O(\|X\|^3).
\]
\end{lemma}

\begin{proof}
Use the Taylor expansion $\exp(X)=I+X+\tfrac12X^2+O(\|X\|^3)$ and $\mathrm{Tr}(X)=0$ for $X\in\mathfrak{su}(3)$.
Since $X$ is anti-Hermitian, $\mathrm{Tr}(X^2)$ is real.
\end{proof}

\begin{proposition}[Quadratic part of $S_W$ at vacuum]\label{prop:quad}
At $U\equiv I$, the action expanded in $A$ has the form
\[
S_W(\exp(A)) = \frac{\beta}{12}\sum_{x\in\Lambda_L}\sum_{\mu<\nu}\| (dA)_{\mu\nu}(x)\|^2 + O(\|A\|^3).
\]
Equivalently, writing $\|dA\|^2:=\sum_{x,\mu<\nu}\|(dA)_{\mu\nu}(x)\|^2$,
\[
S_W(\exp(A)) = \frac{\beta}{12}\|dA\|^2+O(\|A\|^3).
\]
\end{proposition}

\begin{proof}
For each plaquette, the group element $U_{\mu\nu}(x)$ equals $\exp\!\big((dA)_{\mu\nu}(x)+O(\|A\|^2)\big)$.
Insert this into Lemma~\ref{lem:class}.  The $O(\|A\|^2)$ correction in the exponent contributes only $O(\|A\|^3)$ to the action because the quadratic term depends only on the linear part of the exponent.
\end{proof}

\section{The Hessian operator at vacuum}

Let $H_0:=\nabla^2 S_W(I)$ denote the Hessian at the vacuum, viewed as a self-adjoint operator on $\mathfrak{su}(3)^{E_L}$ with respect to the inner product of Document~01.

\begin{definition}[Adjoint $d^\ast$ and the $1$-form Laplacian]
Define $d^\ast$ as the adjoint of $d$ with respect to the lattice inner products on $1$-forms and $2$-forms induced by $\langle\cdot,\cdot\rangle$ on $\mathfrak{su}(3)$.
Define the (coexact) $1$-form Laplacian by
\[
\Delta_1 := d^\ast d.
\]
\end{definition}

\begin{theorem}[Vacuum Hessian equals discrete $1$-form Laplacian]\label{thm:H0}
The Hessian at the vacuum satisfies
\[
H_0 = \frac{\beta}{6}\,\Delta_1
\]
as operators on $\mathfrak{su}(3)^{E_L}$.
\end{theorem}

\begin{proof}
By Proposition~\ref{prop:quad},
\[
S_W(\exp(A)) = \frac{\beta}{12}\langle dA,dA\rangle + O(\|A\|^3) = \frac{\beta}{12}\langle A,\Delta_1 A\rangle + O(\|A\|^3).
\]
By definition of the Hessian, the quadratic form is
\[
S_W(\exp(A)) = S_W(I) + \frac12\langle A,H_0 A\rangle + O(\|A\|^3).
\]
Matching the quadratic terms gives $H_0=\frac{\beta}{6}\Delta_1$.
\end{proof}

\section{Kernel and spectral gap on the physical subspace}

The operator $\Delta_1=d^\ast d$ has the following exact kernel at finite volume:
\begin{itemize}[leftmargin=*]
\item \emph{Exact (gauge) modes:} $A=d\omega$ (discrete gradient of a $0$-form), which satisfy $dA=dd\omega\equiv 0$.
\item \emph{Harmonic (toron) modes:} constant $1$-forms $A_\mu(x)\equiv c_\mu$, which satisfy both $dA\equiv 0$ and $d^\ast A\equiv 0$.
\end{itemize}

Let $P_0$ be the vacuum projector from Document~01 onto $\ker(G_0^\dagger)$ (divergence-free sector).  Let $T$ be the orthogonal projector onto the toron subspace $\mathcal H_{\mathrm{tor}}$.

Define the \emph{strictly physical} subspace at vacuum as
\[
\mathcal H_{\mathrm{phys}}^{\mathrm{strict}} := \mathrm{Im}(P_0)\cap \ker(T) = \ker(G_0^\dagger)\cap \mathcal H_{\mathrm{tor}}^\perp .
\]

\begin{theorem}[Finite-volume physical gap at vacuum]\label{thm:gap}
On $\mathcal H_{\mathrm{phys}}^{\mathrm{strict}}$, the operator $\Delta_1$ has a strictly positive spectral gap.
More precisely, in Fourier space one finds that the nonzero eigenvalues coincide with the scalar lattice Laplacian symbol
\[
\hat\lambda(k) = 4\sum_{\mu=0}^3 \sin^2\!\Big(\frac{k_\mu}{2}\Big),
\qquad k_\mu=\frac{2\pi n_\mu}{L},\ \ n_\mu\in\{0,1,\dots,L-1\},
\]
restricted to $k\ne 0$ and transverse polarizations.
Hence the smallest positive eigenvalue is
\[
\lambda_{\min}(\Delta_1|\_{\mathcal H_{\mathrm{phys}}^{\mathrm{strict}}}) = 4\sin^2\!\Big(\frac{\pi}{L}\Big).
\]
Consequently, by Theorem~\ref{thm:H0},
\[
\lambda_{\min}(H_0|\_{\mathcal H_{\mathrm{phys}}^{\mathrm{strict}}})
=\frac{\beta}{6}\,4\sin^2\!\Big(\frac{\pi}{L}\Big)
\sim \frac{\beta}{6}\Big(\frac{2\pi}{L}\Big)^2
\qquad (L\to\infty).
\]
\end{theorem}

\begin{proof}
This is a standard discrete Fourier diagonalization argument:
the symbol of the forward difference $\Delta_\mu$ is $e^{ik_\mu}-1$, so $|\Delta_\mu|^2=4\sin^2(k_\mu/2)$.
For each nonzero $k$, the longitudinal polarization is in $\mathrm{Im}(G_0)$ and is removed by $P_0$; the remaining transverse polarizations carry eigenvalue $\hat\lambda(k)$.
The smallest nonzero $|k|$ on the periodic box occurs at $k=(2\pi/L,0,0,0)$ up to symmetry.
\end{proof}

\section{Consequences for uniform-in-$L$ convexity claims}

Theorem~\ref{thm:gap} shows that the vacuum physical gap scales like $L^{-2}$.
Therefore, any claim of an \emph{$L$-uniform} positive lower bound for the Wilson Hessian in a fixed neighborhood of $U\equiv I$ is incompatible with the vacuum linearization unless additional structure is introduced (e.g.\ a separate local mass term, a shrinking neighborhood in $L$, or a different notion of ``physical'' projection).

\end{document}
